% %%%%%%%%%%%%%%%%%%%%%%%%%%%%%%%%%%%%%%%%%%%%%%%%%%%%%%%%%%%%%%%%%%%%%%%%%%%%%%%%%
%
%	siglas.tex 
%  -----------
%
%	Neste arquivo você escreve as siglas do seu trabalho na seguinte forma:
%
%	\item[Sigla] O que siginifica a sigla
%
% %%%%%%%%%%%%%%%%%%%%%%%%%%%%%%%%%%%%%%%%%%%%%%%%%%%%%%%%%%%%%%%%%%%%%%%%%%%%%%%%%
%

\item[2D] Espaço bidimensional
\item[3D] Espaço tridimensional
% --- A
\item[ABNT] Associação Brasileira de Normas Técnicas
\item[AFTES] \textit{Association Française des Tunnels et de L'espace Souterrain}
\item[ANSYS] \textit{Analysis System Software}
\item[APDL] \textit{Ansys Parametric Design Language}
\item[AXI] Estado axissimétrico de deformações
% --- B
% --- C
\item[CV-CF] Método Convergência-Confinamento
% --- D
\item[DP] Drucker-Prager
% --- E
\item[E] Comportamento elástico
\item[EP] Comportamento elastoplástico
\item[EPVP] Comportamento elastoplástico-viscoplástico
% --- F
\item[FORTRAN] \textit{Formula Translation}
% --- G
% --- H
% --- I
\item[ITA] \textit{International Tunneling and Underground Space Association}
% --- J
% --- K
% --- L
% --- M
\item[MC] Mohr-Coulomb
% --- N
\item[NATM] \textit{New Austrian Tunneling Method}
\item[NIM] \textit{New Implicit Method}
% --- O
% --- P
% --- Q
% --- R
% --- S
% --- T
% --- U
\item[USERMAT] \textit{User Defined Material Constitutive Model}
% --- V
% --- W
% --- X
% --- Y
% --- Z